
\documentclass{article}
\usepackage{colortbl}
\usepackage{makecell}
\usepackage{multirow}
\usepackage{supertabular}

\begin{document}

\newcounter{utterance}

\centering \large Interaction Transcript for game `eqbench', experiment `eq\_bench', episode 35 with reflexion{-}Qwen3.5{-}9B.
\vspace{24pt}

{ \footnotesize  \setcounter{utterance}{1}
\setlength{\tabcolsep}{0pt}
\begin{supertabular}{c@{$\;$}|p{.15\linewidth}@{}p{.15\linewidth}p{.15\linewidth}p{.15\linewidth}p{.15\linewidth}p{.15\linewidth}}
    \# & \multicolumn{2}{c}{Player} && \multicolumn{2}{c}{Game Master} \\
    \hline

    \theutterance \stepcounter{utterance}  
    & & & \multicolumn{4}{p{0.6\linewidth}}{
        \cellcolor[rgb]{0.9,0.9,0.9}{
            \makecell[{{p{\linewidth}}}]{
                \texttt{\tiny{[P1$\langle$GM]}}
                \texttt{You have attempted this same task before and did not succeed. These are the plans you wrote after those attempts.} \\
\texttt{Trial \#1: New plan: I failed because I assigned high scores to emotions like Relief and Pride that contradict the tragic context of the dialogue. Scott's mother explicitly accuses him of abandonment and compares him to his deceased father, creating a heavy atmosphere of failure and regret. My strategy should have prioritized analyzing the emotional weight of the mother's accusations ("You're just going to leave me here," "I want my son") to identify Guilt as the dominant response. I should have assigned a score of 0 to Relief and Pride, as there is no evidence of positive feelings or relief in the interaction, and reserved the highest score (10) for Guilt, with Resentment being a secondary, lower{-}intensity reaction to the unfairness of the accusation.} \\
\\ 
\\ 
\\ 
\texttt{Your task is to predict the likely emotional responses of a character in this dialogue:} \\
\\ 
\texttt{Deirdre: Scott, have you thought about this? Really thought about it?} \\
\texttt{Scott: Yes, mom. I have. } \\
\texttt{Deirdre: You're just going to leave me here, all alone?} \\
\texttt{Scott: That's not fair, mom. You know it's not like that.} \\
\texttt{Deirdre: Isn't it? Your father left me, and now you're doing the same.} \\
\texttt{Scott: Dad didn't leave, mom. He died. There's a difference.} \\
\texttt{Deirdre: But the end result's the same, isn't it? I'm left alone.} \\
\texttt{Scott: I'm not abandoning you. I'm just trying to live my life.} \\
\texttt{Deirdre: And what about my life, Scott? Who's going to take care of me?} \\
\texttt{Scott: You've got Mrs. Henderson next door, and I'll hire a nurse if need be.} \\
\texttt{Deirdre: I don't want a nurse, Scott. I want my son.} \\
\texttt{{[}End dialogue{]}} \\
\\ 
\texttt{At the end of this dialogue, Scott would feel...} \\
\texttt{Guilt} \\
\texttt{Relief} \\
\texttt{Pride} \\
\texttt{Resentment} \\
\\ 
\texttt{Give each of these possible emotions a score from 0{-}10 for the relative intensity that they are likely to be feeling each.} \\
\\ 
\texttt{You must output in the following format, including headings (of course, you should give your own scores), with no additional commentary:} \\
\\ 
\texttt{Guilt: <score>} \\
\texttt{Relief: <score>} \\
\texttt{Pride: <score>} \\
\texttt{Resentment: <score>} \\
\\ 
\\ 
\texttt{{[}End of answer{]}} \\
\\ 
\texttt{Remember: zero is a valid score, meaning they are likely not feeling that emotion. You must score at least one emotion > 0.} \\
\\ 
\texttt{Your answer:} \\
            }
        }
    }
     \\ \\

    \theutterance \stepcounter{utterance}  
    & \multicolumn{4}{p{0.6\linewidth}}{
        \cellcolor[rgb]{0.561,0.737,0.561}{
            \makecell[{{p{\linewidth}}}]{
                \texttt{\tiny{[P1$\rangle$GM]}}
                \texttt{Guilt: 9} \\
\texttt{Relief: 0} \\
\texttt{Pride: 0} \\
\texttt{Resentment: 2} \\
            }
        }
    }
    & & \\ \\

    \theutterance \stepcounter{utterance}  
    & & & \multicolumn{2}{p{0.3\linewidth}}{
        \cellcolor[rgb]{0.9,0.9,0.9}{
            \makecell[{{p{\linewidth}}}]{
                \texttt{\tiny{[GM$|$GM]}}
                \texttt{Guilt: 9} \\
\texttt{Relief: 0} \\
\texttt{Pride: 0} \\
\texttt{Resentment: 2} \\
            }
        }
    }
    & & \\ \\

    \theutterance \stepcounter{utterance}  
    & & & \multicolumn{2}{p{0.3\linewidth}}{
        \cellcolor[rgb]{0.9,0.9,0.9}{
            \makecell[{{p{\linewidth}}}]{
                \texttt{\tiny{[GM$|$GM]}}
                \texttt{Guilt: 8} \\
\texttt{Relief: 2} \\
\texttt{Pride: 0} \\
\texttt{Resentment: 6} \\
            }
        }
    }
    & & \\ \\

    \theutterance \stepcounter{utterance}  
    & & & \multicolumn{2}{p{0.3\linewidth}}{
        \cellcolor[rgb]{0.9,0.9,0.9}{
            \makecell[{{p{\linewidth}}}]{
                \texttt{\tiny{[GM$|$GM]}}
                \texttt{game\_result = LOSE} \\
            }
        }
    }
    & & \\ \\

    \theutterance \stepcounter{utterance}  
    & & \multicolumn{4}{p{0.6\linewidth}}{
        \cellcolor[rgb]{0.95,0.95,0.95}{
            \makecell[{{p{\linewidth}}}]{
                \texttt{\tiny{[P1$\leftrightarrow$P]}}
                \texttt{Post{-}episode Reflexion (not sent to GM):} \\
\texttt{New plan: I failed because I assigned a non{-}zero score to Resentment (2), which contradicts the submissive and apologetic tone of Scott's responses. Throughout the dialogue, Scott consistently validates his mother's pain ("That's not fair," "I'm not abandoning you") and offers practical solutions, showing no signs of bitterness or anger. My strategy should have strictly adhered to the evidence: Scott is overwhelmed by the weight of his mother's loneliness and his perceived failure to protect her, making Guilt the only dominant emotion. I should have assigned a score of 10 to Guilt and 0 to Relief, Pride, and Resentment, as there is absolutely no textual support for any positive feelings or resentment in his interactions.} \\
            }
        }
    }
    & \\ \\

\end{supertabular}
}

\end{document}
